% Provenance (never printed): see sources/outline_sources.md (Year 2: eccentricity, obliquity and
% precession; insolation computed from spherical astronomy; the Quaternary cycles; carbon dioxide
% and feedbacks). Orbital elements and tabulated insolation of Berger and Loutre (1991) via NOAA
% NCEI (data/derived/bachelor-2-orbit91, fetch_data Book 4 block); periods after Laskar et al.
% (2004) and the NASA Milankovitch page; history after the NASA Earth Observatory feature (2000);
% terminations after Barker and Knorr (2021); the mid-Pleistocene transition after Clark et al.
% (2006) and Koehler and van de Wal (2020). Figures from figdata/bachelor-2/orbital.py; values of
% the exercises and problem computed in the Book 4 scratch calculator calc/ch25.py.
\chapter{El forzamiento orbital y las glaciaciones}\label{ch:b2:orbital-forcing}

Durante el último millón de años, la forma de la órbita terrestre ha cambiado la luz solar recibida por
el planeta entero en un año en menos de dos partes por mil: entre \num{340.00} y
\qty{340.54}{W.m^{-2}} en el techo de la atmósfera. Durante ese mismo millón de años, la insolación media
diaria de mediados de julio a \ang{65} N ha oscilado entre unos 388 y \qty{500}{W.m^{-2}}, en más de
una cuarta parte. Las glaciaciones siguieron la segunda cifra, no la primera. Este capítulo calcula ambas
a partir de los elementos de la órbita terrestre y se pregunta cómo la luz solar de una latitud en una estación puede
marcar el ritmo del crecimiento y del retroceso de casquetes de hielo que cubren continentes.
% ledger: data:b4-orbital, berger91:readme

\begin{recall}
\textbf{De los libros de física}: la elipse, su semieje mayor $a$ y su excentricidad $e$, las leyes de
Kepler. \textbf{Del volumen de primer año}: la eclíptica y la oblicuidad $\varepsilon$, la
declinación del Sol $\delta$, con $\sin\delta = \sin\varepsilon\sin\lambda$ ($\lambda$ es la longitud verdadera del Sol,
contada a lo largo de la eclíptica desde el equinoccio vernal), el año trópico, la insolación media
diaria
\[
  \bar Q = \frac{S_0}{\pi}\left(\frac{a}{r}\right)^{2}\bigl(h_0\sin\varphi\sin\delta +
  \cos\varphi\cos\delta\sin h_0\bigr),\qquad \cos h_0 = -\tan\varphi\tan\delta ,
\]
y su media global y anual, $S_0/(4\sqrt{1 - e^2})$. \textbf{De los volúmenes escolares}: el
ciclo glaciar y el Último Máximo Glacial; el forzamiento radiativo. \textbf{De capítulos anteriores}: la
\omterm{def:b2:precambrian-earth:ice-albedo}{retroalimentación hielo--albedo} (\cref{def:b2:precambrian-earth:ice-albedo}); el registro bentónico apilado, el registro de EPICA
Dome C y los \omterm{def:b2:climate-proxies:age-model}{modelos de edad} (\cref{ch:b2:climate-proxies}).
\end{recall}
% ledger: earth:TSI, earth:obliquity, earth:ecc

\section{La órbita y sus tres elementos}\label{sec:b2:orbital-forcing:elements}

La órbita terrestre es una elipse con el Sol en uno de sus focos. Los planetas, sobre todo Júpiter y
Saturno, la deforman lentamente y la hacen girar en su plano; la atracción del Sol y de la Luna sobre el abultamiento
ecuatorial hace que el eje de rotación se balancee como una peonza. Tres magnitudes describen lo que importa para la
luz solar: cuán alargada es la elipse, cuán inclinado está el eje y en qué punto de la órbita
caen las estaciones.
% ledger: nasa:milank-causes

\begin{definition}[Longitud del perihelio; precesión climática]\label{def:b2:orbital-forcing:elements}
La \emph{longitud del perihelio}\index{longitud del perihelio} $\varpi$ es el ángulo, medido en
el plano de la órbita a partir de la dirección del equinoccio vernal, entre esa dirección y la
dirección del perihelio vista desde el Sol; se cuenta desde el equinoccio móvil, de modo que
cambia tanto con el balanceo del eje como con el giro de la órbita. La \emph{precesión
climática}\index{precesion climatica@precesión climática} es la magnitud $e\sin\varpi$, llamada también índice
de precesión.
\end{definition}
% ledger: berger91:readme

\begin{omfigure}
\centering
\begin{tikzpicture}[scale=0.82, font=\scriptsize, >=Stealth]
\draw[thick, rotate=101.37] (-0.72,0) ellipse (2.4 and 2.2895);
\draw[densely dashed, black!60] (-0.331,1.647) -- (0.615,-3.059);
\draw[->, black!70] (0,0) -- (3.9,0);
\node[anchor=south, align=center] at (3.75,0.06) {$\gamma$, hacia el\\equinoccio vernal};
\draw[->, omDef] (0.75,0) arc[start angle=0, end angle=101.37, radius=0.75];
\node[omDef] at (0.78,0.62) {$\varpi$};
\fill[ocRed] (0,0) circle (0.13);
\node[anchor=north west] at (0.05,-0.08) {Sol};
\foreach \x/\y in {-2.062/0, 0/-3.094, 2.321/0, 0/1.688} {\fill[ocBlue] (\x,\y) circle (0.08);}
\node[anchor=east, align=right] at (-2.15,0) {equinoccio\\de marzo};
\node[anchor=north west, align=left] at (2.38,-0.08) {equinoccio\\de septiembre};
\node[anchor=west] at (0.1,1.92) {solsticio de diciembre};
\node[anchor=east] at (-0.12,-3.25) {solsticio de junio};
\fill[black] (-0.331,1.647) circle (0.06);
\node[anchor=east] at (-0.42,1.72) {perihelio};
\fill[black] (0.615,-3.059) circle (0.06);
\node[anchor=west] at (0.7,-3.15) {afelio};
\draw[->, thick, black!60] (-1.9,-1.75) arc[start angle=220, end angle=245, radius=2.6];
\end{tikzpicture}
\par\smallskip{\small La órbita terrestre vista desde el norte, con la excentricidad exagerada
hasta 0.3 y la longitud actual del perihelio, $\varpi = \ang{101.4}$. La Tierra se desplaza en sentido antihorario
(flecha); los puntos marcan su posición en los equinoccios y los solsticios. El perihelio cae unos
días después del solsticio de diciembre, el afelio unos días después del solsticio de junio. Esquema.}
\end{omfigure}
% ledger: berger91:rows, nasa:milank-detail

Visto desde la Tierra, el Sol se desplaza a lo largo de la eclíptica, y su longitud verdadera $\lambda$ es el
ángulo contado desde el equinoccio vernal: $\lambda = \ang{90}$ en el solsticio de junio, \ang{270} en el solsticio de
diciembre. La Tierra está en el perihelio cuando el Sol se halla en la dirección opuesta a la del perihelio,
en $\lambda = \varpi + \ang{180}$; hoy $\varpi = \ang{101.37}$ y el perihelio cae en
$\lambda = \ang{281.37}$, unos once días después del solsticio de diciembre, hacia el 3 de enero.
% ledger: berger91:rows, nasa:milank-detail

\begin{definition}[Parámetros de Milankovitch]\label{def:b2:orbital-forcing:milankovitch}
Los \emph{parámetros de Milankovitch}\index{parametros de Milankovitch@parámetros de Milankovitch} son las tres magnitudes
orbitales que fijan el reparto de la luz solar sobre la Tierra a lo largo de las estaciones y las
latitudes: la excentricidad $e$, la oblicuidad $\varepsilon$ y la \omterm{def:b2:orbital-forcing:elements}{precesión climática}
$e\sin\varpi$. La \emph{teoría de Milankovitch}\index{teoria de Milankovitch@teoría de Milankovitch} afirma que sus lentas
variaciones marcan el ritmo del crecimiento y del retroceso de los casquetes de hielo cuaternarios, a través de la insolación estival
en las latitudes altas del norte.
\end{definition}
% ledger: nasa:milank-detail, eo00:milankovitch

Cada parámetro varía con sus propios ritmos. La excentricidad tiene periodos principales cercanos a \num{95000},
\num{99000}, \num{124000} y \num{131000} años, que se combinan en un ciclo de unos
\num{100000} años, y un periodo de \num{405000} años, el más estable de todos; durante el último
millón de años osciló entre \num{0.0046} y \num{0.0565}, y hoy, con \num{0.0172}, disminuye
lentamente. La oblicuidad varía con un periodo principal de \num{40978} años, entre \ang{22.1} y
\ang{24.5} durante el último millón de años; con \ang{23.44}, está disminuyendo, pasado su último máximo
hace unos \num{10000} años. El eje completa su balanceo en unos \num{25770} años, mientras que la
órbita gira en unos \num{112000}; juntos desplazan el perihelio a través de las estaciones con
periodos principales cercanos a \num{19000}, \num{22400} y \num{23700} años.
% ledger: la04:eccentricity, la04:obliquity, la04:precession, nasa:milank-detail, data:b4-orbital

\begin{omfigure}
\centering
% figdata/bachelor-2/orbital.py part elements
\begin{tikzpicture}
\begin{groupplot}[group style={group size=1 by 3, vertical sep=0.3cm, xlabels at=edge bottom,
  xticklabels at=edge bottom}, width=0.96\linewidth, height=3.1cm, xmin=0, xmax=1000,
  xtick={0,100,...,1000}, tick label style={font=\scriptsize}, label style={font=\small},
  scaled y ticks=false, yticklabel style={/pgf/number format/fixed,
  /pgf/number format/precision=2}]
\nextgroupplot[ylabel={$e$}, ymin=0, ymax=0.085, ytick={0,0.02,0.04,0.06}]
\addplot[omDef, thin] table[x=age_ka, y=ecc] {figdata/out/bachelor-2-orbital-elements.dat};
\node[font=\scriptsize, anchor=north west] at (rel axis cs:0.01,0.97) {a\enspace excentricidad};
\nextgroupplot[ylabel={$\varepsilon$ (\unit{\degree})}, ymin=21.9, ymax=25.4, ytick={22,23,24}]
\addplot[ocBlue, thin] table[x=age_ka, y=obl] {figdata/out/bachelor-2-orbital-elements.dat};
\node[font=\scriptsize, anchor=north west] at (rel axis cs:0.01,0.97) {b\enspace oblicuidad};
\nextgroupplot[ylabel={$e\sin\varpi$}, ymin=-0.07, ymax=0.12, ytick={-0.05,0,0.05},
  xlabel={edad (miles de años antes de 1950)}]
\addplot[ocRed, thin] table[x=age_ka, y=prec] {figdata/out/bachelor-2-orbital-elements.dat};
\node[font=\scriptsize, anchor=north west] at (rel axis cs:0.01,0.97) {c\enspace precesión
climática};
\end{groupplot}
\end{tikzpicture}
\par\smallskip{\small Los tres \omterm{def:b2:orbital-forcing:milankovitch}{parámetros de Milankovitch} durante el último millón de años: (a) la
excentricidad; (b) la oblicuidad; (c) la \omterm{def:b2:orbital-forcing:elements}{precesión climática}, cuya envolvente es la excentricidad.
Datos: Berger y Loutre (1991), NOAA NCEI Paleoclimatology, consultados el 23 de octubre de 2026.}
\end{omfigure}
% ledger: data:b4-orbit91

La \omterm{def:b2:orbital-forcing:elements}{precesión climática} oscila deprisa, pero solo tanto como se lo permite la excentricidad: cuando la órbita es
casi circular, como hace \num{400000} años, apenas importa dónde cae el perihelio. La
figura lo muestra con claridad: la envolvente del panel (c) es el panel (a).
% ledger: data:b4-orbit91

\section{La luz solar a partir de la órbita}\label{sec:b2:orbital-forcing:insolation}

La insolación media diaria recordada más arriba necesita un ingrediente más: la distancia al Sol en un
día dado.

\begin{proposition}[Distancia y longitud del Sol]\label{prop:b2:orbital-forcing:distance}
Cuando la longitud verdadera del Sol es $\lambda$, la distancia Tierra--Sol es
\[
  \frac{r}{a} = \frac{1 - e^2}{1 - e\cos(\lambda - \varpi)} .
\]
En el solsticio de junio ($\lambda = \ang{90}$), a primer orden en $e$,
$(a/r)^2 \approx 1 - 2e\sin\varpi$; en el solsticio de diciembre, $(a/r)^2 \approx 1 + 2e\sin\varpi$.
\end{proposition}

\begin{proof}
En una elipse con el Sol en un foco, $r/a = (1 - e^2)/(1 + e\cos v)$, donde la anomalía verdadera $v$
es el ángulo, en el Sol, desde el perihelio hasta la Tierra. La Tierra se ve desde el Sol en la
dirección opuesta a aquella en la que el Sol se ve desde la Tierra, así que su longitud vista desde el Sol es
$\lambda + \ang{180}$, y $v = \lambda + \ang{180} - \varpi$, $\cos v = -\cos(\lambda - \varpi)$.
En $\lambda = \ang{90}$, $\cos(\lambda - \varpi) = \sin\varpi$ y $(a/r)^2 = (1 -
e\sin\varpi)^2/(1 - e^2)^2 = 1 - 2e\sin\varpi + O(e^2)$; en $\lambda = \ang{270}$ el signo cambia.
\end{proof}

La \omterm{def:b2:orbital-forcing:elements}{precesión climática} es, por tanto, el índice adecuado: mide cuánto más cerca del Sol, o más
lejos, está la Tierra en el solsticio de junio. Cuando $e\sin\varpi$ es grande y positiva, como hoy y
mucho más hace \num{116000} años, el verano boreal cae cerca del afelio y es débil; cuando es
grande y negativa, como hace unos \num{10000} años, el verano boreal cae cerca del perihelio.
% ledger: berger91:rows

\begin{example}[El Sol de junio a \ang{65} N, hoy]\label{ex:b2:orbital-forcing:today}
Hoy $e = 0.017236$, $\varpi = \ang{101.37}$, $\varepsilon = \ang{23.446}$, y las tablas usan
$S_0 = \qty{1360}{W.m^{-2}}$. En el solsticio de junio $\delta = \varepsilon$; en $\varphi =
\ang{65}$, $\cos h_0 = -\tan\ang{65}\tan\ang{23.446} = -0.9301$, así que $h_0 = \ang{158.44} =
\qty{2.7654}{rad}$: el Sol está sobre el horizonte durante $2 \times 158.44/360 \times 24 = 21.1$~horas. El corchete de
la fórmula de la insolación vale $2.7654 \times 0.9063 \times 0.3979 + 0.4226 \times 0.9174 \times 0.3674 =
1.1397$, y a la distancia media $\bar Q = 432.90 \times 1.1397 = \qty{493.4}{W.m^{-2}}$. El
factor de distancia es $(a/r)^2 = 0.9671$, así que $\bar Q = \qty{477.1}{W.m^{-2}}$.
\end{example}
% ledger: berger91:rows, berger91:readme, data:b4-orbital

La segunda ley de Kepler añade un matiz. Cerca del afelio la Tierra se mueve más despacio: la estación que cae
allí es más débil pero dura más. Los veranos boreales duran hoy unos 4.5 días más que los inviernos
boreales, y las primaveras unos tres días más que los otoños.
% ledger: nasa:milank-detail

\begin{proposition}[Las medias anuales no dependen de la precesión]\label{prop:b2:orbital-forcing:annual}
La energía recibida por unidad de superficie en el techo de la atmósfera entre dos longitudes dadas
del Sol no depende de $\varpi$, y la insolación media anual en cualquier latitud es la de una
órbita circular de radio $a$ multiplicada por $1/\sqrt{1 - e^2}$. En un polo, la media anual es
\[
  \bar Q_{\text{polo}} = \frac{S_0\sin\varepsilon}{\pi\sqrt{1 - e^2}} ,
\]
que crece con la oblicuidad.
\end{proposition}

\begin{proof}
Por la segunda ley de Kepler, el radio barre área al ritmo constante $\pi ab/T$, y el área barrida
mientras la Tierra gira un ángulo $\mathrm d\lambda$ es $\frac12 r^2\,\mathrm d\lambda$, así que $\mathrm dt =
T r^2\,\mathrm d\lambda/(2\pi ab)$, con $b = a\sqrt{1 - e^2}$. La insolación en una latitud y una
longitud del Sol dadas es $(a/r)^2$ veces su valor $Q_a(\lambda)$ a la distancia $a$, de modo que la
energía recibida es $\int Q_a(\lambda)(a/r)^2\,\mathrm dt = \frac{T}{2\pi\sqrt{1 - e^2}}\int
Q_a(\lambda)\,\mathrm d\lambda$: la distancia se ha simplificado, y con ella $\varpi$. A lo largo de todo el
año, $\lambda$ recorre \ang{360} y la media es $\frac{1}{2\pi\sqrt{1-e^2}}\int_0^{2\pi}
Q_a\,\mathrm d\lambda$. En el Polo Norte $Q_a = S_0\sin\delta = S_0\sin\varepsilon\sin\lambda$
para $0 < \lambda < \pi$ y 0 en otro caso; la integral vale $2S_0\sin\varepsilon$.
\end{proof}
% ledger: earth:TSI

La precesión traslada luz solar de una estación a otra, nunca la añade ni la quita a un año; la excentricidad
cambia el total anual en un factor $1/\sqrt{1 - e^2}$, como mucho un \qty{0.16}{\percent} durante el
último millón de años; la oblicuidad traslada luz solar entre latitudes: eleva la media anual en los
polos y la reduce ligeramente en el ecuador. Solo el reparto estacional y latitudinal cambia
mucho, y eso es lo que muestra la figura siguiente.
% ledger: data:b4-orbital, nasa:milank-detail

\begin{omfigure}
\centering
% figdata/bachelor-2/orbital.py part lat
\begin{tikzpicture}
\begin{groupplot}[group style={group size=2 by 1, horizontal sep=1.6cm},
  tick label style={font=\scriptsize}, label style={font=\small}, height=5.4cm,
  width=0.48\linewidth, xmin=-90, xmax=90, xtick={-90,-60,...,90}, xlabel={latitud
  (\unit{\degree})}, ylabel={$\bar Q$ (\unit{W.m^{-2}})}]
\nextgroupplot[ymin=0, ymax=650, title={a\enspace solsticio de junio}, title style={font=\scriptsize},
  legend style={font=\scriptsize, at={(0.03,0.97)}, anchor=north west, draw=none},
  legend cell align=left]
\addplot[black, thick] table[x=lat, y=jun_0] {figdata/out/bachelor-2-orbital-lat.dat};
\addlegendentry{hoy}
\addplot[ocRed, thick, densely dashed] table[x=lat, y=jun_10] {figdata/out/bachelor-2-orbital-lat.dat};
\addlegendentry{hace \num{10000} años}
\nextgroupplot[ymin=100, ymax=500, title={b\enspace media anual}, title style={font=\scriptsize},
  legend style={font=\scriptsize, at={(0.5,0.03)}, anchor=south, draw=none},
  legend cell align=left]
\addplot[ocBlue, thick] table[x=lat, y=ann_lo] {figdata/out/bachelor-2-orbital-lat.dat};
\addlegendentry{$\varepsilon = \ang{22.1}$}
\addplot[ocRed, thick, densely dashed] table[x=lat, y=ann_hi] {figdata/out/bachelor-2-orbital-lat.dat};
\addlegendentry{$\varepsilon = \ang{24.5}$}
\end{groupplot}
\end{tikzpicture}
\par\smallskip{\small (a) Insolación media diaria en el solsticio de junio, hoy y hace \num{10000}
años, cuando el verano boreal caía cerca del perihelio y la oblicuidad era mayor: unos
\qty{50}{W.m^{-2}} más en las latitudes altas del norte. (b) Insolación media anual para las dos
oblicuidades extremas del último millón de años, con la excentricidad actual: una inclinación mayor envía
más luz solar a los polos y algo menos a los trópicos. Calculadas a partir de los elementos orbitales
con $S_0 = \qty{1360}{W.m^{-2}}$.}
\end{omfigure}
% ledger: data:b4-orbital, data:b4-orbit91

\section{El cálculo de la curva de insolación}\label{sec:b2:orbital-forcing:curve}

La curva de insolación que construyó Milankovitch es hoy un cálculo rutinario. Los elementos orbitales del
pasado se calculan integrando los movimientos de los planetas; tabulados cada mil años,
son los datos de partida. Para cada edad, el método siguiente da la insolación en cualquier latitud y estación.

\begin{method}[La insolación del pasado]\label{met:b2:orbital-forcing:curve}
\begin{enumerate}
  \item Léanse $e$, $\varpi$ y $\varepsilon$ para la edad buscada.
  \item Elíjase la estación mediante la longitud verdadera del Sol $\lambda$: \ang{90} para el solsticio de junio;
        la fecha de mediados de julio de las tablas corresponde a $\lambda = \ang{120}$.
  \item Calcúlese $\delta$ a partir de $\sin\delta = \sin\varepsilon\sin\lambda$, y el ángulo horario de la
        puesta de sol $h_0$ en la latitud $\varphi$ (día o noche polar si $|\tan\varphi\tan\delta| > 1$).
  \item Calcúlese el factor de distancia $(a/r)^2$ con la \cref{prop:b2:orbital-forcing:distance}.
  \item Multiplíquese: $\bar Q = (S_0/\pi)(a/r)^2(h_0\sin\varphi\sin\delta +
        \cos\varphi\cos\delta\sin h_0)$.
\end{enumerate}
\end{method}
% ledger: berger91:readme, data:b4-orbital

Aplicado a las 1001 edades del último millón de años, el método reproduce la insolación tabulada de mediados de julio
a \ang{65} N con un margen de \qty{0.012}{W.m^{-2}}, el redondeo de las tablas, y los
valores tabulados de junio en la actualidad con exactitud. Las tablas sitúan el origen del tiempo en 1950.
% ledger: data:b4-orbital, berger91:readme

\begin{omfigure}
\centering
% figdata/bachelor-2/orbital.py part jul65
\begin{tikzpicture}
\begin{groupplot}[group style={group size=1 by 2, vertical sep=0.3cm, xlabels at=edge bottom,
  xticklabels at=edge bottom}, width=0.96\linewidth, xmin=0, xmax=400, xtick={0,50,...,400},
  tick label style={font=\scriptsize}, label style={font=\small}]
\nextgroupplot[height=5.0cm, ymin=370, ymax=590, ylabel={$\bar Q$ a \ang{65} N
  (\unit{W.m^{-2}})}, legend style={font=\scriptsize, at={(0.5,1.02)}, anchor=south, draw=none,
  legend columns=3, /tikz/every even column/.append style={column sep=6pt}},
  legend cell align=left]
\addplot[black!45, line width=2.2pt] table[x=age_ka, y=tab] {figdata/out/bachelor-2-orbital-jul65.dat};
\addlegendentry{mediados de julio, tabulada}
\addplot[ocRed, thin] table[x=age_ka, y=calc] {figdata/out/bachelor-2-orbital-jul65.dat};
\addlegendentry{mediados de julio, calculada}
\addplot[ocBlue, thin, densely dashed] table[x=age_ka, y=jun65] {figdata/out/bachelor-2-orbital-jul65.dat};
\addlegendentry{solsticio de junio, calculada}
\nextgroupplot[height=2.8cm, ymin=-0.02, ymax=0.02, ytick={-0.01,0,0.01},
  scaled y ticks=false, yticklabel style={/pgf/number format/fixed}, ylabel={residuo},
  xlabel={edad (miles de años antes de 1950)}]
\addplot[only marks, mark=*, mark size=0.35pt, black] table[x=age_ka, y=resid]
  {figdata/out/bachelor-2-orbital-jul65.dat};
\end{groupplot}
\end{tikzpicture}
\par\smallskip{\small Arriba: insolación media diaria a \ang{65} N durante los últimos \num{400000} años,
a mediados de julio, tabulada por Berger y Loutre (gris, ancha) y calculada con el método anterior (rojo),
y en el solsticio de junio (azul, discontinua). Abajo: valores de mediados de julio calculados menos tabulados, en
\unit{W.m^{-2}}.}
\end{omfigure}
% ledger: data:b4-orbit91, data:b4-orbital

Destacan dos rasgos. La curva está dominada por la \omterm{def:b2:orbital-forcing:elements}{precesión climática}, con picos separados unos
\num{20000} años cuya altura crece y mengua con la excentricidad; la oblicuidad añade
una modulación más lenta. Y las oscilaciones son grandes: durante los últimos \num{800000} años, el valor de mediados de julio
varió entre 388 y \qty{497}{W.m^{-2}}.
% ledger: data:b4-orbital

\section{La teoría de Milankovitch y los ciclos glaciares}\label{sec:b2:orbital-forcing:theory}

\begin{history}[La insolación a sesenta y cinco grados norte]
Milutin Milankovitch, nacido en 1879, construyó una teoría matemática del clima a partir de las variaciones estacionales y latitudinales de la luz solar. Combinando las tres variaciones orbitales,
calculó la insolación de las latitudes y las temperaturas superficiales correspondientes para
los \num{600000} años anteriores a 1800. Por sugerencia del climatólogo Köppen eligió la insolación
estival a \ang{65} N como magnitud decisiva: los grandes casquetes de hielo habían crecido cerca de esa
latitud, y los veranos frescos funden menos nieve invernal, de modo que la nieve se acumula de un año
a otro. Consideraba la oblicuidad el más importante de los tres parámetros. Durante unos cincuenta años
la teoría fue en gran parte ignorada; Milankovitch murió en 1958. En 1976, Hays, Imbrie y Shackleton
encontraron en sedimentos de mar profundo de los últimos \num{450000} años picos espectrales de variabilidad climática
en torno a \num{23000}, \num{42000} y \num{100000} años, los periodos de la precesión, la oblicuidad y la
excentricidad, y concluyeron que los cambios de la órbita terrestre son la causa fundamental de las
glaciaciones cuaternarias.
\end{history}
% ledger: eo00:milankovitch, nasa:milank-detail, hays76:periods

La figura siguiente repite su prueba con los registros del \cref{ch:b2:climate-proxies}. Compara los
espectros de amplitud, es decir, la intensidad de cada periodo, de la insolación de julio a \ang{65} N y del
registro bentónico apilado, en dos ventanas de \num{700000} años.

\begin{omfigure}
\centering
% figdata/bachelor-2/orbital.py part spectrum
\begin{tikzpicture}
\begin{axis}[width=0.9\linewidth, height=5.6cm, xmin=10, xmax=150, ymin=0, ymax=1.08,
  xlabel={periodo (miles de años)}, ylabel={amplitud relativa},
  xtick={19,23,41,100,150}, tick label style={font=\scriptsize}, label style={font=\small},
  legend style={font=\scriptsize, at={(0.5,1.02)}, anchor=south, draw=none, legend columns=3,
  /tikz/every even column/.append style={column sep=6pt}}, legend cell align=left]
\addplot[black!50, thick] table[x=period, y=amp_ins] {figdata/out/bachelor-2-orbital-spectrum.dat};
\addlegendentry{insolación \ang{65} N, 0--\qty{700}{ka}}
\addplot[ocRed, thick] table[x=period, y=amp_young] {figdata/out/bachelor-2-orbital-spectrum.dat};
\addlegendentry{LR04, 0--\qty{700}{ka}}
\addplot[ocBlue, thick, densely dashed] table[x=period, y=amp_old] {figdata/out/bachelor-2-orbital-spectrum.dat};
\addlegendentry{LR04, \numrange{1250}{1950}~ka}
\end{axis}
\end{tikzpicture}
\par\smallskip{\small Espectros de amplitud, cada uno escalado a su propio máximo, de la insolación de mediados de julio
a \ang{65} N durante los últimos \num{700000} años y del registro bentónico apilado LR04 en la misma ventana
y entre hace \num{1250000} y \num{1950000} años (sin tendencia, con ventana de suavizado). La insolación tiene su
potencia en los periodos de la precesión, entre unos 19 y 24 mil años, y en el periodo de la oblicuidad; el
volumen de hielo reciente tiene su pico cerca de \num{100000} años, donde la insolación casi no tiene nada; el
volumen de hielo más antiguo tiene su pico en \num{41000} años.}
\end{omfigure}
% ledger: data:b4-orbital, data:lr04, hays76:periods

\subsection*{Periodos glaciares, interglaciares y terminaciones}

\begin{definition}[Periodo glaciar; interglaciar; terminación glaciar]\label{def:b2:orbital-forcing:glacials}
En el registro cuaternario, un \emph{periodo glaciar}\index{periodo glaciar} es un intervalo en el que grandes
casquetes de hielo cubren los continentes del norte, registrado por un $\delta^{18}$O bentónico alto y un nivel del mar bajo;
un \emph{interglaciar}\index{interglaciar} es un intervalo con poco hielo continental en el hemisferio
norte fuera de Groenlandia, como el actual. Una \emph{terminación
glaciar}\index{terminacion glaciar@terminación glaciar} es una de las grandes transiciones que cierran un ciclo glaciar,
de condiciones de pleno glaciar a condiciones interglaciares, mucho más rápida que el crecimiento que la precedió.
\end{definition}
% ledger: kohler20:mpt, barker21:terminations, ipcc:ref-periods

Durante el último medio millón de años, el periodo dominante del clima está cerca de \num{100000} años,
aunque el forzamiento en ese periodo, que procede solo de la excentricidad, es despreciable
frente al de la oblicuidad y la precesión. Los ciclos son además asimétricos: el $\delta^{18}$O bentónico
registra una acumulación escalonada de hielo durante unos \num{90000} años, seguida de una desaparición en unos
\num{10000}, un diente de sierra cuyos flancos descendentes son las terminaciones. El último \omterm{def:b2:orbital-forcing:glacials}{interglaciar}, desde
hace unos \num{129000} hasta hace unos \num{116000} años, terminó cuando la insolación estival a \ang{65} N cayó
hasta uno de sus valores más bajos de los últimos \num{800000} años: el valor del solsticio de junio alcanzó su
mínimo, \qty{438.6}{W.m^{-2}}, hace \num{116000} años, y el de mediados de julio su segundo mínimo
más bajo, \qty{392.2}{W.m^{-2}}, hace \num{114000} años. El Último Máximo Glacial llegó unos
\num{95000} años después.
% ledger: barker21:terminations, ipcc:ref-periods, data:b4-orbital, data:b4-orbital-extremes

¿Con qué rapidez puede un casquete de hielo seguir su forzamiento? El modelo más sencillo trata el volumen de hielo $V$ como
si se relajara hacia el volumen $V_{\mathrm{eq}}$ que la insolación del momento sostendría, con un
tiempo de respuesta $\tau$:
\[
  \frac{\mathrm dV}{\mathrm dt} = \frac{V_{\mathrm{eq}}(t) - V}{\tau} .
\]

\begin{proposition}[Una respuesta con retraso]\label{prop:b2:orbital-forcing:lag}
Para un forzamiento $V_{\mathrm{eq}} = A\sin\omega t$, la solución del modelo con retraso tiende a
\[
  V = \frac{A}{\sqrt{1 + \omega^2\tau^2}}\,\sin\bigl(\omega(t - t_{\ell})\bigr),\qquad
  t_{\ell} = \frac{\arctan(\omega\tau)}{\omega} :
\]
el hielo sigue casi por completo un forzamiento lento y solo débilmente un forzamiento rápido, siempre con retraso. La respuesta
a una suma de sinusoides es la suma de las respuestas: el modelo no crea ningún periodo que falte en su
forzamiento.
\end{proposition}

\begin{proof}
Ensáyese $V = B\sin\omega t + C\cos\omega t$: al sustituir, $\omega B\tau = -C$ y $-\omega C\tau = A -
B$, así que $B = A/(1 + \omega^2\tau^2)$, $C = -\omega\tau A/(1 + \omega^2\tau^2)$, una amplitud
$\sqrt{B^2 + C^2} = A/\sqrt{1 + \omega^2\tau^2}$ y un desfase $\arctan(\omega\tau)$. Cualquier otra
solución difiere de esta en un múltiplo de $\mathrm e^{-t/\tau}$, que se extingue. La
ecuación es lineal, así que las soluciones se suman.
\end{proof}

\begin{omfigure}
\centering
% figdata/bachelor-2/orbital.py part lag
\begin{tikzpicture}
\begin{axis}[width=0.96\linewidth, height=5.4cm, xmin=0, xmax=400, ymin=-3.4, ymax=4.2,
  xtick={0,50,...,400}, xlabel={edad (miles de años antes de 1950)},
  ylabel={estandarizado}, tick label style={font=\scriptsize}, label style={font=\small},
  legend style={font=\scriptsize, at={(0.5,1.02)}, anchor=south, draw=none, legend columns=3,
  /tikz/every even column/.append style={column sep=6pt}}, legend cell align=left]
\addplot[black!35, thin] table[x=age_ka, y=forcing] {figdata/out/bachelor-2-orbital-lag.dat};
\addlegendentry{menos la insolación \ang{65} N}
\addplot[ocRed, thick] table[x=age_ka, y=model] {figdata/out/bachelor-2-orbital-lag.dat};
\addlegendentry{modelo con retraso, $\tau = \qty{10}{kyr}$}
\addplot[black, thick] table[x=age_ka, y=lr04] {figdata/out/bachelor-2-orbital-lag.dat};
\addlegendentry{LR04 (hielo hacia arriba)}
\end{axis}
\end{tikzpicture}
\par\smallskip{\small El modelo con retraso forzado por la insolación de mediados de julio a \ang{65} N (menos
sol, más hielo), frente al registro bentónico apilado, los tres estandarizados (media 0, desviación típica
1); los valores positivos significan más hielo. El modelo sigue las oscilaciones de precesión del registro apilado, pero no tiene
ninguno de sus grandes ciclos en diente de sierra. El tiempo de respuesta de \qty{10}{kyr} es una elección del modelo.}
\end{omfigure}
% ledger: data:b4-orbital, data:lr04

Durante los últimos \num{800000} años, el modelo con retraso se correlaciona con el registro apilado con $r = 0.42$, y la
insolación misma solo con $0.23$: el retraso ayuda, pero un casquete de hielo lineal no puede producir ciclos de
\num{100000} años a partir de un forzamiento que no los tiene. Las terminaciones deben implicar un
comportamiento no lineal que amplifique un empuje modesto de la insolación. Los candidatos son retroalimentaciones que actúan sobre todo el clima:
la \omterm{def:b2:precambrian-earth:ice-albedo}{retroalimentación hielo--albedo} de la \cref{def:b2:precambrian-earth:ice-albedo}, ya que un planeta con más
nieve y hielo refleja más luz solar; y el dióxido de carbono, que en los testigos de hielo antárticos sube en
escalas de tiempo de milenios cada vez que se debilita la circulación de vuelco del Atlántico, como ocurre
cuando el agua de fusión de los casquetes de hielo se vierte en el Atlántico Norte. Por encima de un umbral de dióxido
de carbono, las condiciones \omterm{def:b2:orbital-forcing:glacials}{interglaciares} pueden estabilizarse.
% ledger: data:b4-orbital, barker21:terminations, eo00:milankovitch

\begin{definition}[Transición del Pleistoceno medio]\label{def:b2:orbital-forcing:mpt}
La \emph{transición del Pleistoceno medio}\index{transicion del Pleistoceno medio@transición del Pleistoceno medio} es el cambio, desde hace unos
\num{1250000} hasta hace unos \num{700000} años, de ciclos glaciares de unos \num{41000} años y pequeña
amplitud a ciclos cuasiperiódicos de unos \num{100000} años y gran amplitud, con casquetes de hielo más
gruesos, sin ningún cambio correspondiente en el forzamiento orbital.
\end{definition}
% ledger: clark06:mpt, nasa:milank-detail, kohler20:mpt

Antes de la transición, los casquetes de hielo seguían la oblicuidad, como muestra el espectro discontinuo de la figura
anterior; un recuento halla una \omterm{def:b2:orbital-forcing:glacials}{terminación glaciar} en el 67, el 88 y el \qty{52}{\percent} de los
ciclos de oblicuidad del Cuaternario temprano, medio y tardío. La teoría orbital por sí sola no puede explicar el cambio; los candidatos son la dinámica
de los propios casquetes de hielo, quizá tras la eliminación del regolito blando sobre el que se deslizaban,
y un descenso gradual del dióxido de carbono durante los mínimos glaciares. Todavía no hay ninguna explicación
aceptada de forma general.
% ledger: kohler20:mpt, nasa:milank-detail, clark06:mpt

\section{El dióxido de carbono en los ciclos glaciares}\label{sec:b2:orbital-forcing:co2}

Los testigos de hielo antárticos muestran que el dióxido de carbono sigue los ciclos glaciares: entre 188 y
\qty{194}{ppm} en el Último Máximo Glacial, entre 266 y \qty{282}{ppm} en el último \omterm{def:b2:orbital-forcing:glacials}{interglaciar};
el registro compuesto del \cref{ch:b2:climate-proxies} da \qty{189.7}{ppm} para el intervalo de 19 a
\qty{23}{ka} y \qty{279.4}{ppm} para los dos mil años anteriores a la industria. Su forzamiento radiativo
crece casi como el logaritmo de la concentración, \qty{3.93}{W.m^{-2}} por cada duplicación.
% ledger: ipcc:fig2-34, data:b4-proxies-co2, ipcc:log-co2, ipcc:erf-2xco2

\begin{example}[La órbita frente al dióxido de carbono]\label{ex:b2:orbital-forcing:co2}
De \qty{189.7}{ppm} a \qty{279.4}{ppm}, el forzamiento es $3.93\log_2(279.4/189.7) = 3.93 \times
0.559 = \qty{2.2}{W.m^{-2}}$, sobre todo el globo y todo el año. La órbita, durante el último
millón de años, cambió la media global y anual de la insolación como mucho en $340.54 - 340.00 =
\qty{0.54}{W.m^{-2}}$, pero la insolación de mediados de julio a \ang{65} N en unos \qty{109}{W.m^{-2}}.
\end{example}
% ledger: data:b4-proxies-co2, ipcc:erf-2xco2, ipcc:log-co2, data:b4-orbital

La órbita es un marcapasos, no un motor. Redistribuye la luz solar entre estaciones y latitudes
y fija el calendario; los casquetes de hielo, por su albedo y su lentitud, y el dióxido de carbono, por
su forzamiento global, convierten una señal regional en un cambio global del clima.
% ledger: hays76:periods, barker21:terminations, eo00:milankovitch

\section*{Ejercicios}\label{sec:b2:orbital-forcing:exercises}

Datos de los ejercicios, salvo aviso contrario: $S_0 = \qty{1360}{W.m^{-2}}$ (el valor de las
tablas); un año trópico de \num{365.2422} días; los elementos orbitales de las tablas, cuyos valores actuales son
$e = 0.017236$, $\varpi = \ang{101.37}$, $\varepsilon = \ang{23.446}$.
% ledger: berger91:readme, berger91:rows

\begin{exercise}[$\star$]\label{exo:b2:orbital-forcing:1}
Tómese $e = 0.0167$. (a) Calcúlese la razón entre la distancia del afelio, $a(1 + e)$, y la distancia
del perihelio, $a(1 - e)$. (b) Exprésese su diferencia como porcentaje de $a$. (c) ¿En qué factor es
la luz solar en el perihelio más intensa que en el afelio?
\end{exercise}
% ledger: earth:ecc, nasa:milank-detail

\begin{exercise}[$\star$]\label{exo:b2:orbital-forcing:2}
Las tablas dan $e = 0.017236$, $\varpi = \ang{101.37}$ hoy; $e = 0.019971$, $\varpi =
\ang{293.86}$ hace \num{10000} años; $e = 0.019249$, $\varpi = \ang{359.99}$ hace \num{6000} años.
(a) Calcúlese la \omterm{def:b2:orbital-forcing:elements}{precesión climática} para cada edad. (b) Para cada una, dese la longitud del Sol en el
perihelio. (c) ¿En qué edad estaba el perihelio más cerca del solsticio de junio, y en cuál del
equinoccio de septiembre?
\end{exercise}
% ledger: berger91:rows

\begin{exercise}[$\star$]\label{exo:b2:orbital-forcing:3}
El eje de rotación precesa con un periodo de unos \num{25770} años respecto de las estrellas, mientras que
la órbita gira con un periodo de unos \num{112000} años. Supóngase que ambos movimientos desplazan el
perihelio respecto del equinoccio en el mismo sentido, de modo que sus frecuencias se suman. (a) Combínense
las dos frecuencias para hallar el periodo con el que el perihelio
recorre las estaciones. (b) Compárese con los periodos principales de la \omterm{def:b2:orbital-forcing:elements}{precesión climática}, \num{19.0},
\num{22.4} y \num{23.7} mil años.
\end{exercise}
% ledger: nasa:milank-detail, la04:precession

\begin{exercise}[$\star$]\label{exo:b2:orbital-forcing:4}
Con la \cref{prop:b2:orbital-forcing:annual} y $e = 0.0167$, calcúlese la insolación media anual en
el Polo Norte para oblicuidades de \ang{22.1} y \ang{24.5}. ¿En qué porcentaje crece entre
ambas?
\end{exercise}
% ledger: nasa:milank-detail, earth:ecc

\begin{exercise}[$\star\star$]\label{exo:b2:orbital-forcing:5}
(a) Calcúlese exactamente el factor de distancia $(a/r)^2$ en el solsticio de junio hoy y hace \num{116000}
años ($e = 0.043988$, $\varpi = \ang{92.71}$), y compárese con la aproximación $1 -
2e\sin\varpi$. (b) Calcúlense ambos en el solsticio de diciembre. (c) ¿Por qué es peor la aproximación hace \num{116000}
años?
\end{exercise}
% ledger: berger91:rows

\begin{exercise}[$\star\star$]\label{exo:b2:orbital-forcing:6}
Reprodúzcase el \cref{ex:b2:orbital-forcing:today} paso a paso: el ángulo horario de la puesta de sol a \ang{65} N
en el solsticio de junio, la duración del día, la insolación a la distancia media y la
insolación actual.
\end{exercise}
% ledger: berger91:rows

\begin{exercise}[$\star\star$]\label{exo:b2:orbital-forcing:7}
(a) Usando $\mathrm dt = T(r/a)^2\,\mathrm d\lambda/(2\pi\sqrt{1 - e^2})$ y $(r/a)^2 \approx 1 +
2e\cos(\lambda - \varpi)$, demuéstrese que el semestre que va del equinoccio de marzo al de septiembre dura
aproximadamente $T/2 + (2T/\pi)\,e\sin\varpi$. (b) Calcúlese en cuántos días supera hoy al otro semestre.
(c) Compárese con los 4.5 días en que los veranos boreales superan a los inviernos y los tres días en que
las primaveras superan a los otoños.
\end{exercise}
% ledger: nasa:milank-detail, berger91:rows

\begin{exercise}[$\star\star$]\label{exo:b2:orbital-forcing:8}
Durante el último millón de años, la excentricidad osciló entre \num{0.0046} y \num{0.0565}. (a) Calcúlese
la insolación media global y anual $S_0/(4\sqrt{1 - e^2})$ en ambos extremos. (b) ¿Cuánto la cambió la órbita, en
\unit{W.m^{-2}} y en porcentaje? (c) Compárese con el intervalo de la
insolación de mediados de julio a \ang{65} N, de 388 a \qty{497}{W.m^{-2}}.
\end{exercise}
% ledger: data:b4-orbital

\begin{exercise}[$\star\star$]\label{exo:b2:orbital-forcing:9}
Con el modelo con retraso y $\tau = \qty{10}{kyr}$, calcúlense la amplitud relativa y el retraso del
volumen de hielo para un forzamiento sinusoidal de periodos 23, 41 y \num{100} mil años. ¿Qué banda
favorece el casquete de hielo, y por qué?
\end{exercise}

\begin{exercise}[$\star\star\star$]\label{exo:b2:orbital-forcing:10}
La insolación de mediados de julio a \ang{65} N casi no tiene potencia en \num{100000} años, y sin embargo el volumen de hielo
de los últimos \num{700000} años tiene allí la mayor parte de la suya. (a) Demuéstrese que si el forzamiento es $A_1\sin
\omega_1 t + A_2\sin\omega_2 t$, el modelo con retraso responde solo en $\omega_1$ y $\omega_2$. (b) Demuéstrese
que una respuesta que contiene el término $V^2$, por ejemplo $\mathrm dV/\mathrm dt = (V_{\mathrm{eq}} -
V)/\tau + kV^2$, puede crear las frecuencias $\omega_1 - \omega_2$ y $\omega_1 + \omega_2$,
sustituyendo la solución lineal en $V^2$. (c) ¿Con qué periodo baten dos términos de precesión de periodos \num{23.7} y
\num{19.0} mil años? ¿Qué se concluye sobre los ciclos de
\num{100000} años?
\end{exercise}
% ledger: la04:precession, barker21:terminations

\begin{exercise}[$\star\star\star$]\label{exo:b2:orbital-forcing:11}
En el Polo Norte, la energía recibida por metro cuadrado durante el día polar es $E = \int S_0(a/r)^2
\sin\delta\,\mathrm dt$ a lo largo del semestre en que $\delta > 0$. (a) Demuéstrese, con la demostración de la
\cref{prop:b2:orbital-forcing:annual}, que $E = S_0T\sin\varepsilon/(\pi\sqrt{1 - e^2})$. (b) Calcúlese
$E$ en \unit{GJ.m^{-2}} para $\varepsilon = \ang{22.1}$ y \ang{24.5}, con $e = 0.0167$ y $T$ =
365.25 días. (c) La precesión no aparece en $E$. ¿Por qué puede, aun así, cambiar la
fusión de un casquete de hielo?
\end{exercise}
% ledger: nasa:milank-detail

\begin{exercise}[$\star\star\star$]\label{exo:b2:orbital-forcing:12}
(a) Calcúlese el forzamiento radiativo de la subida del dióxido de carbono de \qty{189.7}{ppm} a
\qty{279.4}{ppm}, con \qty{3.93}{W.m^{-2}} por duplicación y una ley logarítmica. (b) Los intervalos
dados para el Último Máximo Glacial y el último \omterm{def:b2:orbital-forcing:glacials}{interglaciar} son de 188 a \qty{194}{ppm} y de 266 a
\qty{282}{ppm}: calcúlense el forzamiento más pequeño y el más grande que permiten. (c) Compárese con el cambio
de la insolación global anual del \cref{exo:b2:orbital-forcing:8} y con la oscilación de la
insolación estival a \ang{65} N. ¿Cuál es el desencadenante y cuál el amplificador?
\end{exercise}
% ledger: data:b4-proxies-co2, ipcc:fig2-34, ipcc:erf-2xco2, ipcc:log-co2, data:b4-orbital

\begin{problem}[{Problema de fin de semana --- La insolación a \ang{65} N}]\label{pb:b2:orbital-forcing:1}
El último \omterm{def:b2:orbital-forcing:glacials}{interglaciar} terminó hace unos \num{116000} años. ¿Cuánta insolación recibía entonces el verano de la
latitud de los grandes casquetes de hielo del norte, comparado con hoy? Úsense los datos de los
ejercicios. Las tablas dan, hace \num{116000} años, $e = 0.0440$, $\varpi = \ang{92.7}$,
$\varepsilon = \ang{22.52}$; hoy, redondeados del mismo modo, $e = 0.0172$, $\varpi = \ang{101.4}$,
$\varepsilon = \ang{23.45}$.

\medskip\noindent\textbf{Parte I --- La órbita entonces y ahora.}
\begin{enumerate}
  \item Calcúlese la \omterm{def:b2:orbital-forcing:elements}{precesión climática} entonces y ahora.
  \item Calcúlense las distancias del perihelio y del afelio, en unidades de $a$, entonces y ahora.
  \item ¿En qué longitud del Sol estaba la Tierra en el perihelio entonces, y ahora? ¿Cerca de qué solsticio?
  \item ¿En qué factor era la luz solar en el perihelio más intensa que en el afelio, entonces y ahora?
  \item ¿Era la oblicuidad mayor o menor que hoy? ¿Qué significa eso para los veranos de las
        latitudes altas?
  \item Calcúlese la insolación media global y anual entonces y ahora. ¿Es significativa la diferencia?
\end{enumerate}

\medskip\noindent\textbf{Parte II --- La distancia en el solsticio de junio.}
\begin{enumerate}[resume]
  \item Demuéstrese a partir de la \cref{prop:b2:orbital-forcing:distance} que en el solsticio de junio $r/a = (1 -
        e^2)/(1 - e\sin\varpi)$.
  \item Calcúlese $r/a$ en el solsticio de junio entonces y ahora.
  \item Calcúlese $(a/r)^2$ entonces y ahora.
  \item Compárese con la aproximación $1 - 2e\sin\varpi$.
  \item ¿En qué porcentaje era la luz solar del solsticio más débil que a la distancia media, entonces y ahora?
\end{enumerate}

\medskip\noindent\textbf{Parte III --- La insolación a \ang{65} N.}
\begin{enumerate}[resume]
  \item Dese la declinación del Sol en el solsticio de junio entonces y ahora.
  \item Calcúlese el ángulo horario de la puesta de sol a \ang{65} N ese día, entonces y ahora.
  \item Dedúzcase la duración del día.
  \item Calcúlese la insolación media diaria a la distancia media.
  \item Calcúlese la insolación media diaria, entonces y ahora.
  \item Las tablas dan, para mediados de julio, 397.60 y \qty{426.76}{W.m^{-2}}. ¿Por qué son más altos los
        valores del solsticio?
  \item ¿Qué pesó más en la diferencia entre entonces y ahora, la oblicuidad o la distancia?
\end{enumerate}

\medskip\noindent\textbf{Parte IV --- La respuesta del casquete de hielo.}
\begin{enumerate}[resume]
  \item En el modelo con retraso con $\tau = \qty{10}{kyr}$, ¿qué fracción de un cambio brusco de
        $V_{\mathrm{eq}}$ ha alcanzado el casquete de hielo al cabo de 5, 10 y \qty{20}{kyr}?
  \item La insolación de junio a \ang{65} N cayó desde un máximo hace unos \num{128000} años hasta el
        mínimo de hace \num{116000} años, medio ciclo de precesión. Para una sinusoide de periodo
        \qty{23}{kyr}, ¿qué fracción del forzamiento sigue el casquete de hielo, y con cuánto retraso?
  \item ¿Por qué no puede un casquete de hielo alcanzar, en un solo ciclo de precesión, el tamaño que una insolación
        estival baja sostendría?
  \item El hielo tardó unos \num{90000} años en acumularse hasta el Último Máximo Glacial y unos
        \num{10000} años en fundirse. ¿Por qué no puede el modelo con retraso producir esa asimetría?
  \item El dióxido de carbono estaba entre 266 y \qty{282}{ppm} en el último \omterm{def:b2:orbital-forcing:glacials}{interglaciar} y entre
        188 y \qty{194}{ppm} en el Último Máximo Glacial. ¿Cómo ayuda a que un cambio regional de la
        insolación se convierta en un cambio global del clima?
  \item En resumen: ¿cuál era la insolación media diaria a \ang{65} N en el solsticio de junio cuando
        terminó el último \omterm{def:b2:orbital-forcing:glacials}{interglaciar}, hace \num{116000} años?
\end{enumerate}
\end{problem}
% ledger: berger91:rows, berger91:readme, data:b4-orbital, ipcc:ref-periods, ipcc:fig2-34, barker21:terminations
